% Generated draft: requires mathematical and source review.
\documentclass[11pt]{article}
\usepackage[margin=25mm]{geometry}
\usepackage{fontspec}
\setmainfont{Latin Modern Roman}
\usepackage{amsmath,amssymb,mathtools,unicode-math}
\setmathfont{Latin Modern Math}
\usepackage{enumitem,needspace,xurl,hyperref}
\setlength{\parindent}{0pt}
\setlength{\parskip}{5pt}
\setlength{\emergencystretch}{3em}
\newcommand{\N}{\mathbb{N}}
\newcommand{\Z}{\mathbb{Z}}
\newcommand{\Q}{\mathbb{Q}}
\newcommand{\R}{\mathbb{R}}
\newcommand{\C}{\mathbb{C}}
\newcommand{\F}{\mathbb{F}}
\newcommand{\A}{\mathbb{A}}
\newcommand{\PP}{\mathbb{P}}
\newcommand{\E}{\mathbb{E}}
\newcommand{\eps}{\varepsilon}
\newcommand{\dd}{\,\mathrm{d}}
\begin{document}
\section*{UnsolvedMath ID: 20002887}
\textbf{Source problem code:} catalogue entry
\subsection{Definitions and mathematical statement}
Let $(M,p)\subset\C^n$ be a generic real-analytic submanifold germ, meaning $T_aM+iT_aM=T_a\C^n$, and let $n\leq N$. Write $B^N=\{w\in\C^N:\sum_j|w_j|^2<1\}$. Suppose $f:M\to\partial B^N$ is continuous and CR and is the restriction of a meromorphic mapping on a neighborhood of $p$ wherever that mapping is defined. Francine Meylan asks when $f$ extends holomorphically to a full neighborhood of $p$, in particular whether minimality of $M$ suffices. Here minimality means that the local CR orbit has full dimension in $M$. The source comments that the assertion holds when $M$ has real codimension $1$, citing [Ch].

The supplied partial reduction can be stated using a reduced common-denominator representation
\[
F=(P_1,\ldots,P_N)/q,
\]
where $P_j,q$ are holomorphic near $p$, $q\not\equiv0$, and no irreducible holomorphic factor of $q$ divides all the $P_j$. Sphere-valuedness gives the real-analytic identity
\[
\sum_{j=1}^N|P_j|^2=|q|^2\quad\text{on }M.
\]
Let $h$ be an irreducible denominator factor and $H=\{h=0\}$. If the trace $M\cap H$ contains a real-analytic submanifold $S$ in the smooth locus of the relevant irreducible hypersurface, generic within $H$ and a holomorphic uniqueness set there, then the aid asserts that every $P_j$ vanishes on $H$. Thus $h$ divides all numerators, contradicting reducedness. Such a factor cannot remain as a nonremovable pole. The uniqueness-set requirement is relative to the complex hypersurface, not just genericity in ambient $\C^n$.

Under minimality, the aid states that the required uniqueness condition follows whenever the trace has an ordinary rank-two zero: a nearby point $a\in M\cap H$ with
\[
\operatorname{rank}_{\R}\dd(h|_M)_a=2,
\]
where $h|_M$ is viewed as a map to $\C\simeq\R^2$. Consequently a reduced counterexample to holomorphic extension with minimal source would necessarily satisfy
\[
\operatorname{rank}_{\R}\dd(h|_M)_a<2
\quad\text{at every nearby }a\in M\cap\{h=0\}
\]
for every irreducible factor $h$ still occurring in its denominator, with irreducible representatives and minimality taken on a sufficiently small neighborhood. This is a necessary restriction on a hypothetical counterexample. It does not remove the remaining factors whose entire zero trace is critical, and therefore does not settle the minimal-source question.
\subsection{Short English statement}
Ask when a continuous sphere-valued CR map with a meromorphic extension is actually holomorphic near the source point. The supplied aid removes pole factors whose trace is a uniqueness set, and under minimality restricts any surviving factor to a wholly critical trace.
\subsection{Sources}
\begin{itemize}
\item Local UnsolvedMath registry, permanent ID $20002887$, source code \texttt{catalogue entry}; canonical \texttt{aim-several-complex-variables-notes.json}, \texttt{notes[44]}.
\item AIM workshop \emph{Complexity of mappings in CR geometry}, Problem 5, \emph{Holomorphic continuation of CR maps}, proposed by Francine Meylan: \url{https://aimath.org/WWN/complexitycrmap/complexitycrmap.pdf}. The source's codimension-one comment cites [Ch]; no full entry is supplied. Original remarks and literature field are empty.
\item Supplied research attempt 1: \texttt{reduction}, problem status at run \texttt{unknown}, current registry status \texttt{partially\_solved}, classification \texttt{PARTIAL-PROGRESS}, rigor \texttt{proved}, novelty confidence low. The divisor-trace criterion is machine-generated aid, not peer reviewed, and requires independent verification before citation.
\end{itemize}
\end{document}
